% Generated draft: requires mathematical and source review.
\documentclass[11pt]{article}
\usepackage[margin=25mm]{geometry}
\usepackage{fontspec}
\setmainfont{Latin Modern Roman}
\usepackage{amsmath,amssymb,mathtools,unicode-math}
\setmathfont{Latin Modern Math}
\usepackage{enumitem,needspace,xurl,hyperref}
\setlength{\parindent}{0pt}
\setlength{\parskip}{5pt}
\setlength{\emergencystretch}{3em}
\newcommand{\N}{\mathbb{N}}
\newcommand{\Z}{\mathbb{Z}}
\newcommand{\Q}{\mathbb{Q}}
\newcommand{\R}{\mathbb{R}}
\newcommand{\C}{\mathbb{C}}
\newcommand{\F}{\mathbb{F}}
\newcommand{\A}{\mathbb{A}}
\newcommand{\PP}{\mathbb{P}}
\newcommand{\E}{\mathbb{E}}
\newcommand{\eps}{\varepsilon}
\newcommand{\dd}{\,\mathrm{d}}
\begin{document}
\section*{UnsolvedMath ID: 30003526}
\textbf{Source problem code:} catalogue entry
\subsection{Definitions and mathematical statement}
Let \(S^2=\{x\in\R^3:\|x\|=1\}\). Under the conventional interpretation using smooth Riemannian metrics, let \(\mathcal G\) be the set of all smooth Riemannian metrics on \(S^2\). For \(g\in\mathcal G\), define
\[
\operatorname{sys}(S^2,g)=\inf\{L_g(\gamma):\gamma\text{ is a nonconstant closed }g\text{-geodesic}\},
\qquad
\rho(S^2,g)=\frac{\operatorname{sys}(S^2,g)^2}{\operatorname{Area}_g(S^2)}.
\]
Here \(L_g(\gamma)\) denotes geodesic length and \(\operatorname{Area}_g(S^2)\) denotes total Riemannian area. The systole in this formulation concerns closed geodesics without a noncontractibility requirement, since every loop on \(S^2\) is contractible.

With these conventions, the question is whether
\[
\sup_{g\in\mathcal G}\rho(S^2,g)=2\sqrt{3}.
\]
The supplied record does not explicitly define \(\rho\), specify metric regularity, or state whether singular metrics are admissible. Thus the displayed formulation is conditional on the stated conventions and requires confirmation against the cited source. The background also mentions contact forms and local optimality of Zoll forms; it supplies no contact-geometric definitions or normalization, so that broader discussion does not specify an additional precise problem here.
\subsection{Short English statement}
Using squared shortest closed-geodesic length divided by total area as the systolic ratio, is its supremum over all smooth Riemannian metrics on the two-sphere equal to \(2\sqrt{3}\)? The ratio convention and admissible metric class require source confirmation.
\subsection{Sources}
\begin{itemize}
\item Supplied local UnsolvedMath registry record, permanent ID \(30003526\), title: Optimal Systolic Ratio of the Two-Sphere; exact source problem code: \texttt{catalogue entry}. The collection is Oberwolfach Reports Open Problems. The proposed year is recorded as 2017; no individual proposer is supplied.
\item Supplied report citation: Dynamische Systeme (2018), Oberwolfach Reports, pp. 1987--2046, DOI 10.4171/owr/2017/32. \url{https://doi.org/10.4171/owr/2017/32}. The background describes Oberwolfach Reports as records of mathematical workshops and their open-problem discussions. The same link is listed among the literature sources as Oberwolfach Report 32/2017.
\item Supplied literature source: Abbondandolo--Benedetti, Local systolic optimality. \url{https://doi.org/10.1007/s00039-023-00624-z}.
\item Literature assessment dated exactly 2026-08-21: the supplied curation labels the problem open, describes the supremum for metrics or contact forms on the two-sphere as unknown, and states that Zoll forms are locally maximizing without determining the global supremum. These are supplied dated curation claims, not independently verified facts. The record expressly requests independent verification of time-sensitive status and cited claims.
\item The registry labels the statement \texttt{corrected\_verified} and reports a manual semantic, typographic, and LaTeX review claiming self-containment and preservation of the target. This is supplied verification metadata; the present draft does not independently establish those claims or resolve the missing conventions.
\end{itemize}
\end{document}
